\subsubsection{Free-Space Scene}
\label{subsubsec:controlled_validation_free_space}

The results in Table \ref{tab:exp1_fs_prop} shows that enabling ISAC reduces the mean path loss from $102.39$ dB to the range $84.39$--$86.82$dB.
The mean path delay remains at $54.10$ ns as the mean distance is same for all the variants.
The differences of performance is more observable in received power as in path loss.
Among the tested ablations, \texttt{isac\_wide} achieves the lowest path loss and the strongest received power, which indicates that the wider beam is beneficial when clutter is minimal.

\begin{table}[t]
    \centering
    \caption{Free-space propagation metrics.}
    \label{tab:exp1_fs_prop}
    \footnotesize
    \begin{tabular}{lccc}
        \hline
        Variant           & Path loss [dB] & Delay [ns] & Rx power [dBm] \\
        \hline
        baseline          & 102.39         & 54.10      & $-56.39$       \\
        isac\_full        & 86.30          & 54.10      & $-40.30$       \\
        isac\_no\_mti     & 85.53          & 54.10      & $-39.53$       \\
        isac\_narrow      & 86.82          & 54.10      & $-40.82$       \\
        isac\_wide        & 84.39          & 54.10      & $-38.39$       \\
        isac\_deep        & 85.54          & 54.10      & $-39.54$       \\
        isac\_no\_tracker & 85.43          & 54.10      & $-39.43$       \\
        \hline
    \end{tabular}
\end{table}

The sensing and tracking results in Table~\ref{tab:exp1_fs_sense} follow a similar trend.
The tracking-enabled ISAC variants achieve detection rates between $0.847$ and $0.903$, while their false-positive rates remain close to $0.018$--$0.019$.
Among them, \texttt{isac\_deep} achieves the highest detection rate $0.903$ and lowest missed-detection rate $0.097$, which shows that additional propagation depth can help to detect better but the gain over the other tracked variants is small.
However, the extra computation necessary for deeper propagation does not worth the gain.
The \texttt{isac\_no\_tracker} variant achieves the highest raw detection rate ($0.970$), but it also increases the false-positive rate to $0.176$.
This means that the tracker is useful because it trades raw sensitivity for cleaner detection.

\begin{table}[t]
    \centering
    \caption{Free-space sensing metrics.}
    \label{tab:exp1_fs_sense}
    \footnotesize
    \begin{tabular}{lccccc}
        \hline
        Variant           & Det.  & FPR   & Miss  & Err. [m] & Cont. \\
        \hline
        baseline          & 0.000 & 0.000 & 1.000 & --       & 0.000 \\
        isac\_full        & 0.847 & 0.019 & 0.153 & 1.148    & 0.550 \\
        isac\_no\_mti     & 0.897 & 0.018 & 0.103 & 1.125    & 0.537 \\
        isac\_narrow      & 0.897 & 0.018 & 0.103 & 1.125    & 0.537 \\
        isac\_wide        & 0.897 & 0.018 & 0.103 & 1.125    & 0.537 \\
        isac\_deep        & 0.903 & 0.018 & 0.097 & 1.126    & 0.510 \\
        isac\_no\_tracker & 0.970 & 0.176 & 0.030 & 1.208    & 0.583 \\
        \hline
    \end{tabular}
\end{table}